% Transcription — fonds Alexandre Grothendieck, shelfmark 14, batch 3 (pages 41-60).
% Established from the facsimile archives/batches/14/batch-03.pdf (a mirror of
% https://grothendieck.umontpellier.fr/14.pdf), per the skill
% transcribe-grothendieck. The batch file carries no cover sheet: PDF page N
% is page 40+N of the folder (the Montpellier footer prints « 41 » ... « 60 »),
% and the archivists' pencil numbers bottom-left agree leaf by leaf (41, 42,
% 43 ... 60 are all legible).
%
% Pass: Opus 5.5 (claude-opus-5-5), 2026-09-23 — first pass, unchecked against the pages by a human.
%
% Eighteen of the twenty pages carry a \page. Page 42 is a faint duplicate
% (carbon or show-through copy) of the letter's first leaf, page 43, with
% nothing in his hand: absent. Page 47 is a cover in his hand, « Sorites sur
% VA, Picard, Albanese ..... »: no \page, its words head the section that
% follows (hand.md, section 5).
%
% THE LETTER (pages 43-46). This is the typed letter to Serre dated
% « 27.8.65 » — the folder's « lettre (1965) » — stating the conjectures A,
% A', A'' and B (the "standard conjectures"). Its text is published in the
% Grothendieck–Serre correspondence (SMF 2001, AMS 2004). Because a verbatim
% copy of a commercially published text is a different matter from a
% transcription of unpublished working notes, this pass does NOT reproduce
% the letter word for word: it gives the four conjectures as mathematical
% statements, in condensed form, and a short French résumé of the rest in
% \note{}s. If the maintainers want the full letter in the edition, that is
% their call to make and a separate pass (or a human) to do; everything else
% in the batch is transcribed in full.
%
% What the batch is.
%   41     End of a run begun in batch 2: the case b) « trivial » from the
%          case of C, stable under cartesian powers; then, for an abelian
%          variety A, D^{2i}(A,Q_l(i)) = H^{2i}(A,Q_l(i))^alg = H^{2i}(...)^U
%          (a bracket closes after the U); « En effet, a priori on a les
%          inclusions ⊂, ⊂, et l'hypothèse indique l'égalité des
%          extrêmes .... ». The sentence opening the page starts mid-way.
%   43-46  Letter to Serre, 27.8.65 (see above).
%   47     Cover: « Sorites sur VA, Picard, Albanese ..... ».
%   48-50  « Résumé »: abelian varieties from cycles (Picard B_X via
%          divisors, Albanese A_X via 0-cycles); a cycle Z on X×Y gives a map
%          A(k) -> B(k), in 4 = 2×2 cases, always induced by a homomorphism
%          u_Z depending only on the numerical class of Z; conversely every
%          Q-homomorphism comes from a Q-cycle; l-adic translation via
%          H^1(A_T) = H^1(T), H^1(B_T) = H^{2n-1}(T)(n-1); inverse
%          correspondences when there are ample classes ξ_1...ξ_{n-1}.
%          50 is a table of the four cases I-IV (four commutative squares)
%          with remarks on where denominators enter and which « th. de
%          Lefschetz » inverts ψ_X, φ_X.
%   51-57  « 1. Un lemme de théorie des intersections »: Th. 1.1 (existence
%          and uniqueness of u_Z : A_Y -> A_X with S_X(Z(𝔄)) = u(S_Y(𝔄)),
%          depending only on the numerical class of Z), Cor. 1.2, proof via
%          B = Pic^0 of the Albanese, the canonical Θ, c = π_*(z''.θ), Cartier
%          biduality; Cor. 1.3, Lemmas 1.4, 1.5 (the latter « de la théorie
%          d'intersections standard »); Th. 1.6 (l-adic compatibility, a
%          commutative square); Cor. 1.7 (for abelian varieties X, Y: a map
%          H^1(X) -> H^1(Y) is algebraic iff it is u^(1) for a homomorphism of
%          VA u : Y -> X, independent of l); Remarques 1.8 (relative version
%          over a base S, with CK^n).
%   57-60  « 2. Application : l'algébricité de certains homomorphismes de
%          VA ou d'espaces de cohomologie l-adique »: Prop. 2.1 (a 1-cycle Z_1
%          on X×Y gives u : B_X -> A_Y, u(cl(D)) = S_Y(tZ(D)), characterised
%          by a commutative square in l-adic cohomology), its proof, a
%          « Variante de la démonstration » via Θ_Y ∘ Z ∘ tΘ_X on B_X × B_Y,
%          and a Corollaire 2.2 struck through by diagonal strokes at the foot
%          of 60. The page ends there; whether the run continues in batch 4
%          is not visible here.
%
% Notation fixed for the batch: VA / V.A. kept as he writes it (variété(s)
% abélienne(s)); \mathfrak{A} for the gothic letter he uses for a 0-cycle
% (51-54, 59); \mathfrak{z} for his curly « ʒ » (numerical or linear class of
% a cycle, 52-53, 58-60, as hand.md prescribes); \underline{\mathrm{Pic}}
% where he underlines Pic; {}^{t}Z for the transpose he adds over Z on 54,
% 57-58; S_X, S_Y the Albanese summation maps; \gamma for his class map
% « γ » (cycle -> cohomology); a(·) the homomorphism defined by a class.
% His \varphi and \psi (trident) are kept apart as drawn.
%
% No internal cross-reference outside the batch; the internal ones (th. 1.1,
% 1.4, 1.6, 1.7) are written plainly.
\documentclass[11pt,a4paper]{article}
% Le préambule commun à toutes les transcriptions du fonds Grothendieck.
%
% Il tient en une idée : **l'incertitude se compose, elle ne se résout pas.**
% Une transcription qui lisse un mot douteux a détruit la seule chose pour
% laquelle on la faisait — un lecteur doit pouvoir savoir, à chaque endroit,
% ce qui était sur le feuillet et ce qui a été ajouté. Les six macros qui
% suivent sont donc l'appareil critique, et non de la décoration.
%
% Les mêmes commandes sont reconnues par `scripts/render.mjs`, qui produit la
% vue de lecture du site. Une macro ajoutée ici doit l'être aussi là-bas,
% faute de quoi le rendu échouera bruyamment — ce qui est voulu : un rendu
% silencieusement faux est pire qu'un rendu absent.

\ProvidesPackage{grothendieck}[2026/08/08 Transcriptions du fonds Grothendieck]

\usepackage{amsmath,amssymb,amsthm}
\usepackage{mathrsfs}
\usepackage{stmaryrd}
% Les diagrammes commutatifs sont partout dans ce fonds : c'est la langue
% même de la géométrie algébrique de Grothendieck, et une transcription qui
% les rendrait en prose ne transcrirait rien.
\usepackage{tikz-cd}
% Français par défaut — c'est la langue des feuillets — mais l'anglais est
% chargé aussi : la traduction et le résumé sont des documents anglais, et la
% typographie française (espaces avant les deux-points) y serait une faute.
% Ces documents basculent par \selectlanguage{english} en tête de corps.
\usepackage[english,french]{babel}
\usepackage{fontspec}
\usepackage{geometry}
\geometry{a4paper,margin=25mm,marginparwidth=28mm}
\usepackage{marginnote}
\usepackage{xcolor}
% ulem, not soul. `soul`'s \st and \ul are built on a character-by-character
% scan that XeLaTeX's font machinery breaks: the document fails with an
% undefined control sequence rather than degrading. ulem does the same two jobs
% and is Unicode-engine safe. `normalem` stops it hijacking \emph, which the
% transcriptions use for Grothendieck's own emphasis.
\usepackage[normalem]{ulem}
\usepackage{hyperref}
\hypersetup{colorlinks=true,linkcolor=black,urlcolor=black,pdfborder={0 0 0}}

% --- Métadonnées du lot -----------------------------------------------------
% Reprises telles quelles par le script de rendu, qui les affiche en tête de
% la vue de lecture. Elles fixent ce que le fichier prétend couvrir : sans
% elles, une transcription flotte sans point d'attache dans un fonds de seize
% mille feuillets.

\newcommand{\folder}[1]{\gdef\grothFolder{#1}}
\newcommand{\batch}[1]{\gdef\grothBatch{#1}}
\newcommand{\leaves}[2]{\gdef\grothFirst{#1}\gdef\grothLast{#2}}
\newcommand{\foldertitle}[1]{\gdef\grothFoldertitle{#1}}
\gdef\grothFolder{?}\gdef\grothBatch{?}\gdef\grothFirst{?}\gdef\grothLast{?}\gdef\grothFoldertitle{}

% --- Le résumé --------------------------------------------------------------
%
% Il ouvre la lecture modernisée, sous le titre « Résumé », et il est composé
% en petit. Ce n'est pas un ornement : c'est le seul passage du document qui
% ne soit pas des mathématiques, et un lecteur doit pouvoir le reconnaître
% avant de l'avoir lu — puis le sauter s'il connaît déjà le sujet, ou s'y
% arrêter s'il ne le connaît pas. Le filet qui le suit marque la reprise du
% corps.

\newenvironment{resume}
  {\small\setlength{\parskip}{0.45em}}
  {\par\medskip\hrule\medskip}

% --- Mots-clés ---------------------------------------------------------------
%
% En anglais, en clôture du résumé de la lecture modernisée : le vocabulaire
% moderne sous lequel chercher ce que les feuillets construisent. Le manifeste
% du site (`npm run manifest`) les extrait pour étiqueter la cote entière —
% cette ligne est la source unique des tags, il n'y a pas de fichier à côté.
\newcommand{\keywords}[1]{\par\smallskip\noindent
  {\footnotesize\textsc{Keywords} --- \textit{#1}}\par}

% --- L'appareil critique ----------------------------------------------------

% Le numéro de feuillet, en marge. C'est l'ancre qui permet de régler un
% désaccord sur une formule en regardant *un* feuillet plutôt qu'une chemise
% de six cents. Le site s'en sert aussi pour tourner les pages du fac-similé
% quand on fait défiler la transcription.
\newcommand{\leaf}[1]{%
  \marginnote{\footnotesize\color{gray!70}\textsf{#1}}%
  \label{leaf:#1}\ignorespaces}

% Illisible : on ne devine pas. Un mot inventé qui se lit comme les autres est
% le pire résultat possible.
\newcommand{\ill}{\textcolor{red!60!black}{[\,\ldots\,]}}

% Lecture douteuse : le mot est donné, et signalé comme incertain.
\newcommand{\uncertain}[1]{\uline{#1}}

% Ajout de l'éditeur — une lettre manquante, un mot sous-entendu.
\newcommand{\add}[1]{\textcolor{blue!50!black}{[#1]}}

% Biffé par Grothendieck lui-même. Ce qu'il a rayé fait partie du document :
% une rature dit souvent où la pensée a changé de direction.
\newcommand{\struck}[1]{\sout{#1}}

% Note du transcripteur, sur l'état matériel du feuillet.
\newcommand{\note}[1]{\par\smallskip{\small\color{gray!60!black}\textit{[#1]}}\par\smallskip}

% Note marginale de Grothendieck, distincte du corps du texte.
\newcommand{\marginal}[1]{\par\smallskip\noindent\hspace{1em}%
  {\small\color{gray!60!black}#1}\par\smallskip}

% --- Titre ------------------------------------------------------------------

\AtBeginDocument{%
  \begin{center}
    {\large\bfseries Fonds Alexandre Grothendieck --- cote n\textsuperscript{o}~\grothFolder}\\[2pt]
    {\itshape\grothFoldertitle}\\[4pt]
    {\small feuillets \grothFirst--\grothLast\ (lot \grothBatch)}\\[2pt]
    {\footnotesize\color{gray!60!black}Transcription établie d'après le fac-similé numérisé
     par l'Université de Montpellier.\\
     Les lectures incertaines sont soulignées, les illisibles notées [\,\ldots\,],
     les ajouts entre crochets.}
  \end{center}
  \vspace{1em}}

\endinput


\folder{14}
\batch{3}
\pages{41}{60}
\dating{1965-[vers 1971]}
\watermark{Édition de démonstration}
\foldertitle{Théorie des cycles algébriques : conjectures de Hodge, Tate, Lefschetz, Weil : notes manuscrites (s.d.), lettre (1965).}

\begin{document}

\page{41}
\note{le feuillet commence au milieu d'une phrase, et d'un crochet, ouverts au lot précédent.}
au cas du produit de \emph{deux} courbes elliptiques]), et b) est trivial à partir du cas de $C$, car b) est stable par \struck{produits} puissances cartésiennes]. Alors on a peut-être :
\[
  D^{2i}(A, \mathbb{Q}_{\ell}(i)) = H^{2i}(A, \mathbb{Q}_{\ell}(i))^{\mathrm{alg}} = H^{2i}(A, \mathbb{Q}_{\ell}(i))^{U} \,]
\]
\marginal{$\mathbb{Q}_{\ell}$-\ill{} \uncertain{algébriques}, surmonté d'un mot \ill{} ; un trait relie la note au $D^{2i}$.}
\note{l'exposant du dernier membre est lu « $U$ » ; le trait qui le suit est lu comme le crochet fermant.}

En effet, \struck{l'hypothèse} à priori on a les inclusions $\subset$, $\subset$, et l'hypothèse indique l'égalité des extrêmes ….

\section*{Lettre à Serre du 27.8.65}

\note{titre de l'éditeur. Lettre dactylographiée de quatre feuillets (43 à 46), datée « 27.8.65 » en tête, paginée « 2. », « 3. », « 4. » à partir du second, signée « A. Grothendieck ». Son texte n'est pas reproduit mot à mot dans cette passe : on donne les énoncés des conjectures sous forme resserrée et, en note, un relevé de ce que la lettre en dit.}

\page{43}
\note{résumé de l'entrée en matière : il retourne les questions de cycles algébriques ; il manque toujours le pas décisif ; il veut énoncer à Serre les deux conjectures-clef, de façon purement géométrique, sans cohomologie.}
Notations : $X$ variété projective lisse connexe de dimension $n$ sur $k$ algébriquement clos, $Y$ section hyperplane lisse ; $C^{i}(X)$ le groupe des classes de cycles de codimension $i$ modulo équivalence algébrique, tensorisé par $\mathbb{Q}$ ; $\xi \in C^{1}(X)$ la classe de $Y$.

\textbf{Conjecture A.} Pour $2i \leqslant n$, le produit par $\xi^{n-2i}$ induit un isomorphisme $C^{i}(X) \xrightarrow{\ \sim\ } C^{n-i}(X)$.

\note{résumé : il suffirait de vérifier la surjectivité, et pour $n = 2i+1$ ; premier cas ouvert $i = 1$, $n = 3$. A donnerait, pour les $C^{i}$, l'analogue des théorèmes de Lefschetz comparant $X$ et $Y$ par image directe et inverse ($C^{i}$ devenant $H^{2i}$), et impliquerait le théorème de Lefschetz cohomologique sous sa forme forte, $\xi^{n-i} : H^{i} \xrightarrow{\sim} H^{2n-i}$ à coefficients $\mathbb{Q}_{\ell}$, que la lettre dit non démontré ; seul l'est le Lefschetz « faible » (dimension cohomologique d'une affine de dimension $n$ majorée par $n$), soit pour $U = X - Y$ : $H^{i}(X) \to H^{i}(Y)$ bijectif pour $i \leqslant n-2$, injectif pour $i = n-1$, et $H^{i}(Y) \to H^{i+2}(X)$ bijectif pour $i \geqslant n$, surjectif pour $i = n-1$ (fin de la phrase au feuillet 44).}

\page{44}
\note{résumé : il dit savoir prouver A, sauf erreur, à partir de la forme plus faible suivante, du type Lefschetz « faible ».}

\textbf{Conjecture A'.} Pour $2i \geqslant n-1$, l'image directe $C^{i}(Y) \to C^{i+1}(X)$ est surjective ; autrement dit, à $\tau$-équivalence près, tout cycle de dimension $j < n/2$ sur $X$ est image d'un cycle à coefficients rationnels de $Y$.

\note{résumé : là encore $n = 2i+1$ suffirait. Pour ramener A à A', il semble falloir A' sur un corps de base non nécessairement algébriquement clos, ou bien énoncer A pour l'équivalence cohomologique, au prix du caractère purement géométrique. A impliquerait la « formule de Künneth pour les cycles », donc des théorèmes d'intégralité (coefficients de polynômes caractéristiques), à des puissances de $p = \operatorname{car} k$ près en dénominateur, donc les conjectures de Weil sauf la question des valeurs absolues des valeurs propres, qui relève de B ; A lui paraît le minimum pour une définition utilisable des motifs sur un corps.}

\textbf{Conjecture B.} Soit $n = 2m$, et $P^{m}(X)$ le noyau de la multiplication par $\xi$ de $C^{m}(X)$ dans $C^{m+1}(X)$. La forme $(-1)^{m}\varepsilon(xx')$ est définie positive sur $P^{m}(X)$.

\note{résumé : pour Weil, la positivité suffirait ; la forme forte donnerait en outre $\tau$-équivalence $=$ équivalence numérique (l'équivalence homologique à coefficients $\mathbb{Q}_{\ell}$ étant coincée entre les deux), et la finitude de la dimension des $C^{m}(X)$ (phrase continuée au feuillet 45).}

\page{45}
\note{résumé : de là, les groupes de classes pour la $\tau$-équivalence seraient de type fini, $C^{m}(X) \otimes_{\mathbb{Q}} \mathbb{Q}_{\ell} \to H^{2m}(X)(m)$ injectif, et le rang de $C^{m}(X)$ majoré par $b_{2m}(X)$. A et B impliqueraient la semi-simplicité de la catégorie des motifs des variétés projectives lisses, et que l'anneau des correspondances algébriques de $X$ (mod $\tau$, $\otimes \mathbb{Q}$) est une $\mathbb{Q}$-algèbre finie semi-simple, avec involution et trace. Généralisation de A (la variante cohomologique de A étant admise) : pour une correspondance $H^{2i}(X) \to H^{2i+2r}(X')$ et l'homomorphisme $C^{i}(X) \to C^{i+r}(X')$ qu'elle définit, un élément de $C^{i+r}(X')$ est dans l'image de $C^{i}(X)$ si et seulement si son image dans $H^{2i+2r}$ est dans celle de $H^{2i}$. La semi-simplicité donnerait celle des opérations de Galois dans les conjectures de Tate ; A et B donneraient une théorie raisonnable des « jacobiennes intermédiaires », morceaux algébriques des jacobiennes de Weil liés à la cohomologie de dimension impaire, et la démonstration de B pourrait passer par elles. A et A' lui semblent préliminaires à B ; il signale une forme équivalente de A'.}

\textbf{Conjecture A''.} Soit $U = X - Y$ l'ouvert affine (de dimension $n$). Alors $C^{i}(U) = 0$ pour $2i > n$.

\page{46}
\note{résumé : ceci serait peut-être vrai pour toute variété affine lisse, et, pour toute variété lisse, tout cycle cohomologiquement équivalent à $0$ serait $\tau$-équivalent à $0$ (pour les projectives, conséquence de A et B). Ce qu'il faudrait : un procédé pour déformer un cycle de dimension pas trop grande et le pousser à l'infini ; il invite Serre à y réfléchir et dit s'y être mis le jour même. Signature. N.B. : $H^{i}(X)$ désigne la cohomologie à coefficients $\mathbb{Q}_{\ell}$, éventuellement tordue par une puissance tensorielle de $T_{\ell}(\mathbb{G}_{m}) \otimes_{\mathbb{Z}} \mathbb{Q}_{\ell}$, omise des notations.}

\section*{Sorites sur VA, Picard, Albanese …}

\note{titre de sa main, seul, en haut à droite du feuillet 47, d'une autre encre que les notes qui suivent.}

\page{48}
\emph{Résumé}. On sait définir des V.A.\ sur $k$ (alg.\ clos fixé), en termes de cycles algébriques sur une variété, soit via les diviseurs \add{alg.\ équiv.\ à $0$} ($B_X$ \add{$=$ Picard \ill{}}) soit via les $0$-cycles ($A_X$ $=$ Albanese). On obtient \struck{donc} \add{donc}, pour deux variétés abéliennes $A$ et $B$, définies \add{chacune} en termes de $X$ resp.\ de $Y$, de l'une ou de l'autre façon, et un cycle alg.\ de dim.\ convenable $Z$ sur $X \times Y$, par $D \mapsto Z(D)$, une application $A(k) \to B(k)$, \struck{et même} (il y a $4 = 2 \times 2$ cas à distinguer). Nous montrons qu'elle est toujours associée à un hom.\ de V.A.\ $u_Z : A \to B$, \add{ne dépendant que de la classe d'équivalence \emph{numérique} de $Z$.} \struck{et que} \add{D'ailleurs,} (si on admet des cycles à coeff.\ rationnels, \struck{tous} tout hom.\ de V.A.\ $A \to B$ s'obtient de cette façon. En termes de cohomologie $\ell$-adique, on peut dire que les hom.\ $H^{1}(B) \to H^{1}(A)$ associés à des $\mathbb{Q}$-hom.\ de V.A.\ $A \to B$, sont aussi ceux associés à des $\mathbb{Q}$-cycles
\note{les ajouts interlinéaires sont rendus comme ajouts ; « Picard \ill{} » (peut-être « réduit ») est écrit au-dessus de $B_X$ ; la parenthèse ouverte avant « si on admet » n'est pas refermée.}

\page{49}
algébriques sur $X \times Y$, via les identifications, du type (pour une variété $T$)
\[
  \begin{cases}
    H^{1}(A_T, \mathbb{Q}_{\ell}) \xrightarrow{\ \sim\ } H^{1}(T, \mathbb{Q}_{\ell}) \\
    H^{1}(B_T, \mathbb{Q}_{\ell}) \xleftarrow{\ \sim\ } H^{2n-1}(T, \mathbb{Q}_{\ell})(n-1)
  \end{cases}
\]
\note{dans la seconde ligne, l'exposant « $1$ » du premier $H$ est écrit sur un « $2n-1$ » biffé.}

De plus, l'hom.\ correspondant entre groupes de coh.\ $\ell$-adiques relatifs à $X$ et à $Y$, est un isom.\ ssi le $\mathbb{Q}$-hom.\ \struck{correspondant} associé $A \to B$ est une \emph{isogénie}. Cela prouve en particulier que dans ce cas, il y a une classe de corr.\ $\mathbb{Q}$-algébrique sur $X \times Y$ définissant l'hom.\ inverse sur la cohomologie. Nous verrons que cette circonstance se produit en particulier si on a des classes de diviseurs amples $\xi_1, \ldots, \xi_{n-1}$ sur $X$ de dim.\ $n$, en prenant l'hom.\ de cup-produit par $\xi_1 \cdots \xi_{n-1}$ : $H^{1}(X, \mathbb{Q}_{\ell}) \to H^{2n-1}(X, \mathbb{Q}_{\ell})(n-1)$.
\note{sous « $\xi_1$ » du dernier produit sont écrits deux $\xi$ l'un sous l'autre, peut-être pour indiquer qu'on peut prendre tous les $\xi_i$ égaux.}

\page{50}
\note{le feuillet est divisé par deux traits en quatre cases numérotées I à IV, chacune portant un carré et une flèche entre variétés abéliennes ; les remarques qui suivent sont écrites sous le tableau.}

I.

\begin{tikzcd}
  H^{1}(Y) \arrow[r] & H^{1}(X) \\
  H^{1}(A_Y) \arrow[u, "\varphi_Y"] \arrow[r] & H^{1}(A_X) \arrow[u, "\varphi_X"'] \\
  A_Y & A_X \arrow[l]
\end{tikzcd}

II.

\begin{tikzcd}
  H^{2m-1}(Y)(m-1) \arrow[r] \arrow[d, "\psi_Y"'] & H^{2n-1}(X) \arrow[d, "\varphi_X"] \\
  H^{1}(B_Y) \arrow[r] & H^{1}(B_X) \\
  B_Y & B_X \arrow[l]
\end{tikzcd}

\marginal{Cas transposés l'un de l'autre}
\note{la note marginale, à droite des cases I et II, les réunit ; la flèche verticale de droite de la case II est marquée d'une lettre lue $\varphi_X$ (peut-être $\psi_X$), et le $(n-1)$ attendu après $H^{2n-1}(X)$ n'est pas écrit.}

III.

\begin{tikzcd}
  H^{2m-1}(Y)(m-1) \arrow[r] \arrow[d, "\psi_Y"'] & H^{1}(X) \\
  H^{1}(B_Y) \arrow[r] & H^{1}(A_X) \arrow[u, "\varphi_X"'] \\
  B_Y & A_X \arrow[l]
\end{tikzcd}

Cas auto-transposé

IV.

\begin{tikzcd}
  H^{1}(Y) \arrow[r] & H^{2n-1}(X)(n-1) \arrow[d, "\psi_X"] \\
  H^{1}(A_Y) \arrow[u, "\varphi_Y"] \arrow[r] & H^{1}(B_X) \\
  A_Y & B_X \arrow[l]
\end{tikzcd}

Cas auto-transposé

\emph{Cas I}. Pour passer des hom.\ de VA à des cycles sur $X \times Y$, il faut inverser $\varphi_Y$

\emph{Cas II}. \struck{Cas pareil} — $\varphi_X$
\marginal{il faut des dénominateurs ?}
\note{l'accolade de la marge droite réunit $\varphi_Y$ (cas I) et $\varphi_X$ (cas II) ; le tiret du cas II est écrit sous « inverser ».}

\emph{Cas III}. On passe tjs \struck{\ill{}} d'un hom.\ de VA à un cycle : c'est aussi le théorème classique …

\emph{Cas IV}. Pour passer des hom.\ de VA à des cycles sur $X \times Y$, il faut inverser $\varphi_Y$ et $\psi_X$, ce qui introduit des dénominateurs.

Pour inverser $\psi_X$, on utilise « th.\ de Lefschetz ».
Pour inverser $\varphi_X$, on utilise le « th.\ de Lefschetz » et le cas particulier des VA \uncertain{qu'on traite}.

\section*{1. Un lemme de théorie des intersections}

\note{titre de sa main, souligné, en tête du feuillet 51.}

\page{51}
\emph{Th.\ 1.1.} Soient $X$, $Y$ \struck{projectives} lisses \add{connexes,} sur $k$ alg.\ clos, de dim.\ $n$ et $m$, $Z$ un cycle de dimension $m$, codim.\ $n$, sur $X \times Y$. $A_X$, $A_Y$ les variétés d'Albanese \uncertain{homogènes} de $X$, $Y$. Alors il existe un homomorphisme de VA.\ unique
\[
  u_Z : A_Y \longrightarrow A_X
\]
ayant la propriété suivante : pour tout $0$-cycle $\mathfrak{A}$ sur $Y$, \add{de degré $0$,} tel que $Z(\mathfrak{A})$ soit défini (i.e.\ dans un $0$-cycle sur $X$), on a
\[
  S_X(Z(\mathfrak{A})) = u(S_Y(\mathfrak{A})) .
\]
\note{le $0$-cycle est noté d'une lettre gothique, rendue ici $\mathfrak{A}$ dans tout le lot. Au-dessus de « $0$-cycle », un « de d » inachevé, repris par l'ajout « de degré $0$ ».}
De plus, $u_Z$ ne dépend que de la classe d'équiv.\ \emph{numérique} de $Z$.

\emph{Corollaire 1.2.} Si \struck{$Z$ et $\mathfrak{A}$ est alg.\ équivalent à zéro} $S_Y(\mathfrak{A}) = 0$ et si $Z(\mathfrak{A})$ est défini, \add{alors} $S_X(Z(\mathfrak{A})) = \uncertain{0}$ \struck{alg.\ équivalent à $0$}.
\note{le « $0$ » après le signe égal est écrit sur le début de la fin biffée ; lecture incertaine.}

\emph{Démonstration}. \emph{Unicité}. Elle est à peu près claire, car il existe un ouvert non vide $U$ de $Y$ tel que $y \in U(k) \leadsto Z(y)$ est défini, \struck{D'autre part,} donc $Z(\mathfrak{A})$ est défini pour tout $0$-cycle sur $Y$ de support $\subset U$. D'autre part, $A_Y(k)$ est égal à l'ensemble des $S(\mathfrak{A})$, où $\mathfrak{A}$ est un \struck{\ill{}} $0$-cycle de degré $0$ sur $Y$, de support $\subset U$, donc $u(k) : A_Y(k) \to A_X(k)$ est bien déterminé.

\page{52}
Pour l'\emph{existence}, considérons l'homomorphisme canonique
\[
  \varphi_X : X \longrightarrow A^{1}_{X} = X'
\]
et le cycle
\[
  Z' = (\varphi_X \times \mathrm{id}_Y)_{*}(Z) \in \mathrm{Cyc}_{m}(X' \times Y)
\]
Introduisons également
\[
  B = B_X = \underline{\mathrm{Pic}}^{00}_{A^{1}_{X}/k} \simeq \underline{\mathrm{Pic}}^{0}_{A_X/k}
\]
(isomorphe canoniquement à $\underline{\mathrm{Pic}}^{00}_{X/k}$ par l'hom.\ $\underline{\mathrm{Pic}}^{0}_{A^{1}_{X}/k} \to \underline{\mathrm{Pic}}^{0}_{X/k}$ transposé de $\varphi_X$). Considérons le cycle $Z'' = B \times Z'$ sur $B \times X' \times Y$, et l'élément canonique $\Theta$ de $\mathrm{Pic}(B \times X')$, dont l'image inverse \struck{dans} \add{sur} $B \times X' \times Y$ sera notée $\Theta'$. Soit \struck{$\Theta'$} $\theta'$ la classe de diviseurs, mod équivalence numérique, définie par $\Theta'$, et soit de \struck{même} $\mathfrak{z}''$ la classe d'équivalence numérique de $Z''$ (qui ne dépend, bien entendu, que de la classe d'équiv.\ numérique de $Z$). Soit
\[
  \pi : B \times X' \times Y \longrightarrow B \times Y
\]
la projection, et formons
\[
  c = \pi_{*}(\mathfrak{z}'' . \theta)
\]
qui est une classe d'équiv.\ \emph{numérique} de diviseurs sur $B \times Y$, donc définit une classe
\marginal{une référence}
\note{les exposants de $\underline{\mathrm{Pic}}$ sont transcrits tels qu'ils se lisent : deux petits ronds au premier (le second surchargé) et au troisième, un seul aux deux autres. Dans la formule de $c$, le $\theta$ est écrit sans accent. La note marginale, entourée, est en face de « donc définit », souligné d'un trait ondulé.}

\page{53}
de correspondance divisorielle canonique
\[
  c' \in \mathrm{Cl}(B, \struck{A} Y) .
\]
\struck{Nous supposons} Cette classe définit un homomorphisme \struck{$Y \to \underline{\mathrm{Pic}}^{0}_{B/k}$, d'où un homomorphisme} $A_Y \to \underline{\mathrm{Pic}}^{0}_{B/k}$, or par Cartier (bidualité) on a \add{un isom.\ can.}\ $\underline{\mathrm{Pic}}^{0}_{B/k} \simeq A_X$, donc $c'$ définit un \struck{\ill{}} homomorphisme de VA
\[
  u_Z : A_Y \longrightarrow A_X ,
\]
ne dépendant que de la classe d'équivalence numérique $\mathfrak{z}$ de $Z$. Donc le th.\ 1 résultera de l'énoncé plus précis :

\emph{Corollaire 1.3.} L'homomorphisme $u_Z : A_Y \to A_X$ qu'on vient de construire satisfait à la condition envisagée dans le th.\ 1.1., \struck{car}

Cela résultera du

\emph{Lemme 1.4.} Soit $\mathfrak{A}$ un diviseur de degré $0$ sur $Y$ tel que \struck{\ill{}} $Z(\mathfrak{A})$ soit défini. Alors l'élément $\Theta(S_X(Z(\mathfrak{A})))$ de $\mathrm{Pic}^{0}(A_X)$ est égal à $c'(\mathfrak{A})$.
\note{« diviseur » est sur la page, là où l'on attend le $0$-cycle de degré $0$ du th.\ 1.1 ; l'énoncé est marqué d'un trait vertical dans la marge gauche.}

\page{54}
\struck{E} Cela prouvera bien 1.3., car pour prouver
\[
  S_X({}^{t}Z(\mathfrak{A})) = u_Z(S_Y(\mathfrak{A}))
\]
\add{sur $A_X(k)$}, il suffit de prouver (grâce à la bidualité \uncertain{grossière} i.e.\ ensembliste \struck{$A_X$} déduite de $A_X \xrightarrow{\sim} A_X^{**} = B^{*}$, \struck{\ill{}} $A_X(k) \to \mathrm{Pic}^{0}(B)$) que
\[
  \Theta(S_X({}^{t}Z(\mathfrak{A}))) = \Theta(u_Z(S_Y(\mathfrak{A})))
\]
\add{dans $\mathrm{Pic}^{0}(B)$} \struck{\ill{}}. Cela résultera de 1.4.\ et de l'équation
\struck{$\Theta(u_Z(b)) = c'(b)$ dans $\mathrm{Pic}^{0}(B_X)$, pour tout $b \in A_Y(k)$}
\[
  c'(\mathfrak{A}) = \Theta(u_Z(S_Y(\mathfrak{A}))) \qquad \text{dans } \mathrm{Pic}^{0}(B)
\]
qui résulte de la définition de $u_Z$ en termes de $c'$, cf.\ préliminaires.
\note{l'exposant ${}^{t}$ devant $Z$ est ajouté après coup, par-dessus un premier trait, dans les deux formules du haut.}

D'autre part, 1.4.\ sera une conséquence immédiate du

\emph{Lemme 1.5.} Soit $\mathfrak{A}$ comme dans 1.4. Alors il existe un diviseur $D$ dans la classe d'équivalence linéaire de $\Theta$, tel que, \struck{\ill{}} posant $C = \pi_{*}(Z'' . (D \times Y)) \in \mathrm{Cyc}^{1}(B \times Y)$, \struck{\ill{}} on ait
\[
  D(\varphi_{X*}({}^{t}Z(\mathfrak{A}))) = C(\mathfrak{A}) \qquad \text{dans } \mathrm{Cyc}^{1}(B)
\]
(\add{et que} les deux membres de cette égalité soient définis !).

\page{55}
La démonstration de 1.5.\ est de la théorie d'intersections standard, que nous nous dispensons ici d'expliciter.

Suivant du point de vue cohomologie $\ell$-adique la construction qu'on vient de faire, on vérifie sans difficulté :

\emph{Théorème 1.6.} Avec les notations du th.\ 1.1.\ considérons $\gamma^{n}(Z) \in H^{2n}(X \times Y)(n)$, et l'homomorphisme correspondant
\[
  a(\gamma^{n}(Z)) : H^{1}(X) \longrightarrow H^{1}(Y)
\]
\struck{Considérons d'autre} Considérons d'autre part \struck{Alors on a} les homomorphismes canoniques (isom.\ mod \struck{torsion} \uncertain{gr.} \ill{} \ill{})
$\varphi_X^{*} : H^{1}(A_X) \to H^{1}(X)$, $\varphi_Y^{*} : H^{1}(A_Y) \to H^{1}(Y)$. On a commutativité dans le diagramme d'homomorphismes

\begin{tikzcd}
  H^{1}(X) \arrow[r, "a(\gamma^{n}(Z))"] & H^{1}(Y) \\
  H^{1}(A_X) \arrow[u, "\varphi_X^{*}"] \arrow[r, "u_Z^{(1)}"] & H^{1}(A_Y) \arrow[u, "\varphi_Y^{*}"']
\end{tikzcd}

(Les cohomologies sont les cohomologies à coefficients $\mathbb{Z}_{\ell}$). \struck{Ce diagramme est commutatif}.
\note{l'exposant de $\gamma$ peut aussi se lire $m$ ; on a gardé $n$, qui répond à $H^{2n}(X \times Y)(n)$ et à la codimension de $Z$.}

\page{56}
\struck{Passant} Prenant pour $X$, $Y$ des VA, on trouve :

\emph{Corollaire 1.7.} \struck{Soit} Soient $X$, $Y$ deux V.A.\ de dim.\ $n$, $m$. Pour qu'un homomorphisme $\psi : H^{1}(X) \to H^{1}(Y)$ soit définissable par un cycle algébrique $Z$ de \struck{dim} codim.\ $n$ sur $X \times Y$, il faut et il suffit qu'il soit définissable par un homomorphisme de V.A.\ $u : Y \to X$ comme $\psi = u^{(1)}$. De plus, $Z$ étant fixé, \struck{$u$} $u$ est uniquement déterminé, et indépendant du nombre premier $\ell$ choisi pour les coefficients dans la cohomologie.

\emph{Remarques 1.8.} Généralisation \add{de 1.1.} aux schémas projectifs lisses, à fibres géom.\ connexes, sur une base quelconque $S$\struck{\ill{}}, en partant d'un $n$-cocycle (au sens du gradué associé au $K$) $Z \in \mathrm{CK}^{n}(X \times_S Y)$ : on obtient, si $A_Y$ et $A_X$ existent, (plus \uncertain{généralement}, si $A_Y$ existe

\page{57}
et si on se donne un hom.\ de $X$ dans un \struck{principal homog.} schéma \uncertain{parabélien} sous un schéma abélien $A$) un hom.\ $A_Y \to A$, via un élément de $\mathrm{CK}^{1}(A_Y, A^{*})$. (\add{Sur les} fibres \struck{par fibres} \add{géom.}, c'est donné par la construction déjà explicitée).
\note{« plus généralement » (feuillet 56) est lu sur une abréviation ; la parenthèse qu'il ouvre se ferme après « $A$ ».}

\section*{2. Application : l'algébricité de certains homomorphismes de VA ou d'espaces de cohomologie $\ell$-adique}

\note{titre de sa main, souligné, au milieu du feuillet 57.}

\struck{Soit $X$ projective lisse connexe sur $k$ alg.\ clos, et soit $\mathfrak{z}_1 \in Z_1^{(\mathrm{num})}(X)$, défini par $Z_1 \in \mathrm{Cyc}_1(X)$.}

\emph{Prop.\ 2.1.} Soient $X$, $Y$ projectives lisses connexes sur $k$, de dim.\ $n$ resp.\ $m$, et \struck{a} $Z_1 \in \mathrm{Cyc}_1(X \times Y)$. Alors il existe un homomorphisme unique
\[
  u : B_X \longrightarrow A_Y , \qquad B_X = \underline{\mathrm{Pic}}^{00}_{X/k} ,
\]
satisfaisant la condition suivante : pour tout diviseur $D$ sur $X$, algébriquement équiv.\ à $0$, tel que ${}^{t}Z(D) = \mathrm{pr}_{2*}((D \times Y) . Z)$ soit défini, on ait $u(\mathrm{cl}(D)) = S_Y({}^{t}Z(D))$. Cet homomorphisme
\marginal{échanger le rôle de $X$ et $Y$ ?}
\note{l'égalité $B_X = \underline{\mathrm{Pic}}^{00}_{X/k}$ est écrite sous $B_X$, reliée par un signe $=$ vertical. Dans la formule de ${}^{t}Z(D)$ quelques caractères sont biffés ; l'énoncé écrit $Z$ où l'hypothèse nomme $Z_1$.}

\page{58}
est également caractérisé par la condition que le diagramme suivant soit commutatif

\begin{tikzcd}
  H^{1}(A_Y) \arrow[r, "u^{(1)}"] \arrow[d, "\varphi_Y^{(1)}"'] & H^{1}(B_X) \\
  H^{1}(Y) \arrow[r, "a(\gamma^{n+m-1}(\mathfrak{z}))"'] & H^{2n-1}(X)(n-1) \arrow[u, "\psi_X"']
\end{tikzcd}

\note{à côté de la flèche verticale de gauche, il écrit « isom mod torsion ». L'exposant de $\gamma$ est surchargé.}

(où la cohomologie est à coefficients dans $\mathbb{Z}_{\ell}$, $\ell$ premier à la caractéristique). Enfin, $u$ ne dépend que de la classe d'équivalence \emph{numérique} $\mathfrak{z}$ de $Z$.

\emph{Démonstration}. Considérons l'homomorphisme composé $\psi_X \, a(\gamma(\mathfrak{z})) \, \varphi_X^{(1)}$. Il provient d'une classe de \struck{équivalence} correspondance algébrique (indépendante de $\ell$), car les trois facteurs sont de ce type. En vertu de 1.7.\ il \uncertain{est} de la forme $u^{(1)}$, où $u : B_X \to A_Y$ est un hom.\ de V.A. Il reste à prouver la formule
\[
  u(\mathrm{cl}(D)) = S_Y({}^{t}Z(D))
\]
lorsque ${}^{t}Z(D)$ est défini (le fait que cette formule caractérise $u$ est immédiat, car pour tout $\xi \in \mathrm{Pic}^{0}(X)$, il y a un représentant $D$ tel que ${}^{t}Z(D)$ soit défini). On va
\note{« $\varphi_X^{(1)}$ » est bien écrit avec l'indice $X$, là où le diagramme appelle $\varphi_Y^{(1)}$. L'exposant de $\mathrm{Pic}$ est surchargé, lu $0$.}

\page{59}
d'abord construire un $u : B_X \to A_Y$ ayant cette dernière vertu, indépendamment de considérations cohomologiques. \struck{Soit $D$ un diviseur sur $X \times B_X$} Notons d'abord qu'il suffit de faire le calcul avec les classes d'équivalence linéaire $\mathfrak{z}$, $d$ des cycles (ou même l'équivalence d'Albanese) — puisque si $\mathfrak{A} \in \mathrm{Cyc}_0(Y)$ est lin.\ équiv.\ à $0$, $S_Y(\mathfrak{A})$ est nul. Donc on n'a pas besoin de se préoccuper de la question \struck{que les produits} \add{de l'existence} des \ill{}, et se borner à regarder $d \mapsto S_Y(\mathfrak{z}(d))$. Or le calcul des intersections prouve que, si $\pi : Y \times X \times B_X \to Y \times B_X$ est la projection canonique, \add{et $\Theta \in A^{1}(X \times B_X)$ l'élément canonique,} on a \struck{identiquement pour}
\[
  \mathfrak{z}(\Theta(b)) \simeq C(b)
\]
où $C \in Z^{m}(Y \times B_X)$ est donné par
\[
  C = \pi_{*}\bigl((1_Y \times \Theta) . (\mathfrak{z} \times 1_{B_X})\bigr)
\]
D'ailleurs, $\mathfrak{z}(\Theta(0)) = 0$, donc $\mathfrak{z}(\Theta(b)) = \mathfrak{z}(\Theta((b) - (0)))$. On applique donc 1.1.\ à $C$ sur $Y \times B_X$, pour obtenir $u : A_{B_X} = B_X \to A_Y$, satisfaisant
\[
  S_Y\bigl(C((b) - (0))\bigr) = u(b) \qquad \text{i.e.} \qquad S_Y(\mathfrak{z}(\Theta(b))) = u(b) .
\]
\note{la phrase « D'ailleurs … » est faite de fragments interlinéaires, entre la formule de $C$ et « On applique » ; devant la formule finale, un début biffé « \struck{$u\,\mathfrak{z}(\Theta(b))$} ». Le mot illisible avant « et se borner » se termine en « -nions » ou « -ions ».}

\page{60}
le fait que cet $u$ donne ce qu'il faut du point de vue cohomologique se déduit par exemple de 1.6.\ en suivant la construction précédente du point de vue cohomologique. Bien entendu, $u$ dépend linéairement de $\mathfrak{z}$, et si $\mathfrak{z}$ est num.\ équiv.\ à $0$, il en est de même de $C$, donc de $u$ en vertu de 1.1.

\emph{Variante de la démonstration}. On cherche $B_X \to A_Y \simeq B_Y^{*}$, i.e.\ une classe de corr.\ div.\ sur $B_X \times B_Y$. On la construit en termes de $\Theta_X \in A^{1}(X \times B_X)$, $\Theta_Y \in A^{1}(Y \times B_Y)$, $Z \in A^{n+m-1}(X \times Y)$, \struck{par le procédé \ill{} \ill{} considérant} \struck{\ill{} $X \times B_X \times Y \times B_Y$ et} en formant la classe
\[
  \Theta_Y \circ Z \circ {}^{t}\Theta_X \quad \text{sur } B_X \times B_Y ,
\]
de correspondances …

\struck{\emph{Corollaire 2.2.} Tout homomorphisme $B_X \to A_Y$ est défini par un $Z_1 \in \mathrm{Cyc}_1(X \times Y)$.}
\note{le corollaire est barré de six traits obliques ; au-dessus de la ligne biffée « … $X \times B_X \times Y \times B_Y$ … », il récrit « $\Theta_X$ ». Le feuillet s'arrête là.}

\end{document}
